\chapter{Fundamentos}
\label{cap:fundamentos}

% \epigrafe{Texto de la cita.}{Autor}

\section{Introducción}

% Introducción del capítulo.

\section{Distribución espacial}
\index{distribucion@distribución!espacial}

\begin{definicion}{Distribución de partículas}{distribucion-particulas}
Sea un sistema de $N$ partículas cuyas posiciones
vienen dadas por $x_i(t)\in\R^d$.
Definimos la distribución espacial normalizada como
\[
\rho(x,t)=\frac{1}{N}\sum_{i=1}^{N}\delta(x-x_i(t)).
\]
\end{definicion}

\begin{ejemplo}{Sistema de dos partículas}{dos-particulas}
Consideremos dos partículas situadas en $a$ y $b$.
La distribución correspondiente es
\[
\rho(x)=\frac12\delta(x-a)+\frac12\delta(x-b).
\]
Su normalización se verifica integrando sobre todo el espacio.
\end{ejemplo}

\begin{teorema}{Normalización}{normalizacion}
La distribución espacial de la Definición~\ref{def:distribucion-particulas}
satisface
\[
\int_{\R^d}\rho(x,t)\,\dif x=1.
\]
\end{teorema}

\begin{demostracion}
Sustituyendo la definición de la distribución,
\[
\int_{\R^d}\rho(x,t)\,\dif x
=\frac1N\sum_{i=1}^N\int_{\R^d}\delta(x-x_i(t))\,\dif x.
\]
Por la propiedad de normalización de la delta de Dirac,
cada integral vale la unidad. Por tanto,
\[
\int_{\R^d}\rho(x,t)\,\dif x=\frac1N\sum_{i=1}^N 1=1.
\]
\end{demostracion}

\begin{observacion}
La distribución no representa necesariamente una densidad ordinaria
respecto a la medida de Lebesgue. Debe interpretarse en sentido
distribucional.
\end{observacion}
